\textit{Proof.}
Put \(s=\rho^2\), and abbreviate \(q_{\sqrt{s}}(A;S)\) and \(v^*(\sqrt{s};S)\) by \(q_s(A)\) and \(v(s)\), respectively. For a unit vector \(u\in H_n\), put \(t(u)=u^\top P_Su\). We first prove
\[
\sup_{0\leq s\leq1}\{q_s(A)-v(s)\}
=\sup_{\substack{u\in H_n\\\|u\|_2=1}}
\{u^\top Au-v(t(u))\}.
\tag{1}
\]
For the lower bound, fix \(u\) and set \(s=t(u)\). Then \(u\) is feasible for \(q_s(A)\), so \(q_s(A)\geq u^\top Au\). For the upper bound, compactness of the spherical shell supplies a maximizer \(u_s\) of \(q_s(A)\). Since \(t(u_s)\geq s\), and since \(v\) is nonincreasing as the adversary class shrinks, we have
\[
q_s(A)-v(s)
=u_s^\top Au_s-v(s)
\leq u_s^\top Au_s-v(t(u_s)).
\]
Taking suprema proves (1).

Enumerate one representative from each balanced sign-pair. By the covariance-hull definition, every \(B\in\mathcal C_n\) can be written as \(A(b)\) with \(b\in\Delta_m\). The corrected scalar dual in \(\text{thm:cone-lift-dual}\), followed by joint minimization over \(b\) and the multiplier, gives, for every \(t\in[0,1]\),
\[
v(t)=
\inf_{\substack{b\in\Delta_m,\ \eta\geq0,\ \ell\in\mathbb R\\
\ell I_{n-1}-Q_H^\top\{A(b)+\eta P_S\}Q_H\succeq0}}
\{\ell-\eta t\}.
\tag{2}
\]
The matrix inequality is equivalent to
\(\ell\geq\lambda_{\max}(A(b)+\eta P_S\mid H_n)\). Formula (2) is an infimum, including at \(t=1\); it need not have a finite-multiplier minimizer there. Substitution of (2) into (1), and the identity that the negative of an infimum is the supremum of the negatives, yield
\[
\Delta_{\mathrm{reg}}(A(a);S)
=\sup_{\substack{u\in H_n,\ \|u\|_2=1;\ b\in\Delta_m,\ \eta\geq0,\ \ell\in\mathbb R\\
\ell I_{n-1}-Q_H^\top\{A(b)+\eta P_S\}Q_H\succeq0}}
\left[u^\top A(a)u+\eta u^\top P_Su-\ell\right].
\tag{3}
\]
This is equality of values; no inner maximizer over finite \(\eta\) is asserted. Minimizing (3) over \(a\in\Delta_m\) and applying \(\text{def:regret-object}\) proves the displayed robust representation. Introducing an epigraph variable \(r\) gives exactly the stated semi-infinite epigraph constraints.

We now make the closure and finite computation claim precise. Set
\[
h=(1+\eta)^{-1},
\qquad
\xi=\ell-\eta.
\]
For \(0<h\leq1\), the matrix constraint in (3), after multiplication by the positive scalar \(h\), becomes
\[
\{1-h+h\xi\}I_{n-1}
-Q_H^\top\{hA(b)+(1-h)P_S\}Q_H\succeq0.
\tag{4}
\]
Let \(Q_S\) be an orthonormal basis matrix for \(S\). The genuine \(h=0\) boundary stratum is
\[
\xi I_{\dim S}-Q_S^\top A(b)Q_S\succeq0.
\tag{5}
\]
To verify (5), write
\[
g_\eta(B)=\lambda_{\max}(B+\eta P_S\mid H_n)-\eta,
\qquad
\kappa_B=\lambda_{\max}(P_SBP_S\mid S).
\]
Every \(B\in\mathcal C_n\) is positive semidefinite and satisfies \(\|B\|_{\mathrm{op}}\leq\operatorname{tr}(B)=n\). Decomposing a unit vector into its \(S\) and \(H_n\cap S^\perp\) components shows, uniformly in \(B\), that for \(\eta>n\),
\[
0\leq g_\eta(B)-\kappa_B
\leq\max_{\delta\geq0}
\{2n\delta-(\eta-n)\delta^2\}
=\frac{n^2}{\eta-n}.
\tag{6}
\]
Thus \(g_\eta(B)\to\kappa_B\) uniformly over the covariance hull. Moreover, \(0\leq g_\eta(B)\leq n\). Since the bracket in (3) decreases with \(\xi\), restricting \(0\leq\xi\leq n\) does not change its supremum.

In the new coordinates the bracket is
\[
u^\top A(a)u+\frac{1-h}{h}\{t(u)-1\}-\xi.
\tag{7}
\]
At \(h=0\), extend (7) to \(u^\top A(a)u-\xi\) when \(u\in S\), and to \(-\infty\) otherwise. Equations (5)--(6) show that this is the value-exact upper-semicontinuous closure of the finite-multiplier carrier. In particular, a value reached only at \(h=0\) is a limit value of (3), not a falsely asserted finite-\(\eta\) maximizer.

The simplexes, the unit sphere, \(h\in[0,1]\), and \(\xi\in[0,n]\) are finite-dimensional compact sets. On \(h>0\), the principal-minor conditions for (4), the sphere equation, the simplex equations and inequalities, and the epigraph inequality obtained from (7) after multiplication by \(h\) are finitely many polynomial sign conditions. The separate boundary stratum (5) is also described by finitely many principal-minor inequalities. Recursively isolating real roots of the coefficient, discriminant, and resultant projection polynomials partitions these conditions into finitely many sign-invariant spectral cells. On each cell the active simplex face is fixed, and every active finite spectral branch is a root of
\[
\det\!\left(\ell I_{n-1}
-Q_H^\top\{A(b)+\eta P_S\}Q_H\right)=0.
\]
Hence the cell boundaries, active faces, generalized eigenvalue crossings, and the compressed boundary (5) give the claimed finite breakpoint description. Exact real-algebraic input permits finite real-algebraic optimization; for a numerical score basis, certified subdivision combined with semidefinite upper and lower bounds computes the value and an \(\varepsilon\)-optimal solution. Every coefficient is obtained from \(n\), the score basis, and the enumerated balanced assignments.

It remains to prove outer attainment and the design-based statement. We first establish continuity from the primitives. If \(S=H_n\), then \(q_s(A)\) is independent of \(s\). Otherwise put \(T=H_n\cap S^\perp\). A unit vector of exact score energy \(t\in[0,1]\) has the form
\[
u=\sqrt t\,x+\sqrt{1-t}\,y,
\]
where \(x\) and \(y\) range over the unit spheres of \(S\) and \(T\). Its quadratic form is
\[
t x^\top Ax+(1-t)y^\top Ay
+2\sqrt{t(1-t)}\,x^\top Ay.
\tag{8}
\]
This is jointly continuous in \((t,A,x,y)\) on a compact set. Parameterizing \(t=s+(1-s)r\), with \(r\in[0,1]\), shows that \((s,A)\mapsto q_s(A)\), the maximum of (8) over a fixed compact parameter set, is continuous. Compact minimization over \(A\in\mathcal C_n\) then makes \(v(s)\) continuous. Equation (1) expresses \(\Delta_{\mathrm{reg}}(A;S)\) as the maximum of a jointly continuous function over the unit sphere, so it is continuous in \(A\). Because \(\mathcal C_n\) is compact, there is an outer minimizer \(A_{\mathrm{reg}}\). By \(\text{lem:balanced-covariance-realizability}\), symmetrizing convex weights for \(A_{\mathrm{reg}}\) produces a design \(d_{\mathrm{reg}}\in\mathcal D_n\).

The exact oracle saddle theorem translates the normalized game into causal mean squared error. For every \(d\in\mathcal D_n\),
\[
\begin{aligned}
&\sup_{\rho\in[0,1]}
\left\{
\sup_{\mu\in\mathcal Y(\rho,M;S)}
\mathbb E_d[(\widehat\tau_d-\tau_n)^2]
-V^*(\rho;S)
\right\}\\
&\qquad=\frac{4M}{n}\Delta_{\mathrm{reg}}(A_d;S)
\geq R^*(n,S).
\end{aligned}
\tag{9}
\]
For the design realizing \(A_{\mathrm{reg}}\), equality holds by the definition of \(R^*(n,S)\). This proves outer attainment and the matching converse over the entire class \(\mathcal D_n\).

Finally, at \(\rho=0\),
\[
q_0(A;S)=\lambda_{\max}(A\mid H_n).
\]
Every \(A\in\mathcal C_n\) is positive semidefinite, annihilates \(\mathbf 1_n\), and has trace \(n\). Therefore
\[
q_0(A;S)\geq\frac{n}{n-1}=c_n.
\]
Complete randomization has covariance \(c_nP_H\) and attains equality, so it is a \(\rho=0\) oracle. Its support is all of \(\mathcal Z_n\). Every symmetrized cut-hull representation of \(A_{\mathrm{reg}}\) uses assignments in \(\mathcal Z_n\), and its support is consequently contained in the support of this single oracle design, hence also in the union of supports of finitely many oracle designs. The causal estimand throughout is the fixed SATE \(\tau_n\); the quantity \(R^*(n,S)\) is the exact design-stage regret for its unbiased estimator. \(\square\)